\documentclass[12pt]{report}
\usepackage[utf8]{inputenc}
\usepackage[T1]{fontenc}
\usepackage{geometry}
\usepackage{titlesec}
\usepackage{enumitem}
\geometry{margin=1in}
\setlength{\parindent}{0pt}
\setlength{\parskip}{0.6em}
\setlength{\emergencystretch}{3em}
\setlist[description]{style=nextline,leftmargin=0pt,labelindent=0pt,font=\bfseries,itemsep=1em}
\titleformat{\chapter}{\normalfont\Large\bfseries\centering}{}{0pt}{}
\titleformat{\section}{\normalfont\large\bfseries}{}{0pt}{}

\begin{document}

\begin{titlepage}
    \centering
    \vspace*{\fill}
    {\Huge\textbf{Summary of Nobel Memorial Prize in Economic Sciences Laureates}}\\[1.5cm]
    {\Large Contributions of the 1969--2024 Winners}\\[2cm]
    \vspace*{\fill}
\end{titlepage}

\tableofcontents
\chapter*{Introduction}
This report briefly describes the main contribution recognized by each Nobel Memorial Prize in Economic Sciences award from 1969 through 2024. Each entry states the main work of the laureate or laureates in clear terms.

{\let\\\relax
\begin{description}

\item[1969 --- Ragnar Frisch and Jan Tinbergen]
Ragnar Frisch helped establish the use of statistics to test ideas about the economy.\\
He developed methods for studying economic relationships and changes over time.\\
Jan Tinbergen built early large models for forecasting and testing government policies.\\
Their work established quantitative testing as an important part of economics. It also created methods still used in economic forecasting and policy evaluation.

\item[1970 --- Paul Samuelson]
Paul Samuelson used mathematics to connect several areas of economic study.\\
He studied consumer choice, public goods, trade, capital, and economic growth.\\
His work also improved the analysis of national income and economic policy.\\
His textbook became an important introduction to modern economics. The book helped standardize the teaching of economic theory in universities.

\item[1971 --- Simon Kuznets]
Simon Kuznets developed methods for measuring national income and economic growth.\\
He collected historical data on output, income, and productivity.\\
His research examined changes in population, cities, industry, and income inequality.\\
The statistical systems he helped develop are still used in economic policy. They also made long-term comparisons between countries more reliable.

\item[1972 --- John Hicks and Kenneth Arrow]
John Hicks developed tools for studying consumer choice and the balance between buyers and sellers.\\
His model explained how spending, interest rates, and the supply of money affect one another.\\
Kenneth Arrow studied how prices can balance buying and selling, and how uncertainty and limited information affect decisions.\\
Together they strengthened the mathematical study of markets and public well-being. Their methods remain important in the study of markets and public policy.

\item[1973 --- Wassily Leontief]
Wassily Leontief developed a method for showing how industries depend on one another.\\
His tables show how one industry's output is used by other industries.\\
He applied this method to planning, trade, employment, and environmental issues.\\
Input--output models are still used to study supply chains and policy effects. They provide a systematic way to trace how a change in one sector affects others.

\item[1974 --- Gunnar Myrdal and Friedrich Hayek]
Gunnar Myrdal analyzed cumulative causation, showing how initial advantages or disadvantages can reinforce themselves.\\
His work on money, cycles, development, and racial inequality combined economics with social analysis.\\
Friedrich Hayek studied monetary disturbances, capital structure, and the limits of centralized knowledge.\\
Their contrasting approaches broadened the study of fluctuations, institutions, and economic coordination. Their work also influenced later research on development and the role of knowledge in markets.

\item[1975 --- Leonid Kantorovich and Tjalling Koopmans]
Leonid Kantorovich applied linear programming and mathematical optimization to resource allocation.\\
He showed how shadow prices could guide the efficient use of scarce inputs.\\
Tjalling Koopmans developed activity analysis and rigorous criteria for productive efficiency.\\
Their work created practical methods for finding the best use of limited resources. These methods are used in both public and private organizations.

\item[1976 --- Milton Friedman]
Milton Friedman changed the study of household spending by showing that people consider their expected long-term income.\\
His monetary history argued that changes in the money supply strongly influence inflation and output.\\
He argued that increasing inflation cannot permanently keep unemployment lower.\\
His empirical and policy work reshaped modern macroeconomic debates. It also influenced the design of monetary policy and the study of inflation.

\item[1977 --- Bertil Ohlin and James Meade]
Bertil Ohlin explained international trade through differences in countries' labor, land, and machines.\\
His analysis clarified specialization, trade distribution, and the effects of protection.\\
James Meade advanced the theory of trade, balance-of-payments adjustment, and optimal economic policy.\\
Together they deepened understanding of how international exchange affects national welfare. Their theories remain basic references in international economics.

\item[1978 --- Herbert Simon]
Herbert Simon showed that decision-makers have limited information and time.\\
He explained how people use simple rules and choose acceptable options.\\
He applied these ideas to firms, organizations, public administration, and computers.\\
His work connected economics with psychology, management, and computer science. It provided a more realistic account of how organizations and individuals make decisions.

\item[1979 --- Theodore Schultz and Arthur Lewis]
Theodore Schultz showed that education, health, and skills increase a person's ability to work and earn.\\
He used this insight to explain agricultural productivity, migration, and development.\\
Arthur Lewis developed a model showing how workers can move from traditional farming into modern industry.\\
Their work placed labor, human capabilities, and structural transformation at the center of development economics. These ideas continue to guide research on poverty and economic growth.

\item[1980 --- Lawrence Klein]
Lawrence Klein created large computer-based models of national economies.\\
He combined economic theory with estimated equations to forecast output, employment, and prices.\\
His models allowed governments and firms to test alternative fiscal and monetary policies.\\
He helped establish empirical macroeconomics as a practical policy tool. His models were used to compare possible economic policies before they were adopted.

\item[1981 --- James Tobin]
James Tobin analyzed how financial markets influence spending, investment, and the wider economy.\\
His work explained how investors divide their money among safer and riskier investments.\\
His measure of company value explained how share prices can affect firms' decisions to invest.\\
He also contributed to liquidity preference, monetary policy, and financial-market regulation. His work linked financial decisions with the wider performance of the economy.

\item[1982 --- George Stigler]
George Stigler studied how firms are organized and how they compete.\\
He explained how information is costly to obtain and how that cost shapes market outcomes.\\
His work on regulation emphasized capture by organized interests rather than purely public motives.\\
He advanced the empirical study of market structure, competition, and government intervention. His analysis remains influential in industrial organization and regulatory economics.

\item[1983 --- Gérard Debreu]
Gérard Debreu gave the study of market balance a precise mathematical form.\\
He established conditions under which prices coordinate consumer and producer choices.\\
His work explained when markets can reach a stable balance, even when choices and technology differ.\\
Debreu's rigorous methods became a model for modern theoretical economics. His work showed how economic conclusions can be derived from clearly stated assumptions.

\item[1984 --- Richard Stone]
Richard Stone developed the modern system of national economic accounts.\\
He defined consistent measures of production, income, expenditure, saving, and capital.\\
His accounting framework made economic activity comparable across sectors and countries.\\
It remains the statistical basis for national income estimates and macroeconomic policy. The framework also supports comparisons of economic performance across countries.

\item[1985 --- Franco Modigliani]
Franco Modigliani developed the life-cycle hypothesis of saving and consumption.\\
He showed how households smooth consumption over their working lives and retirement.\\
His work on corporate finance examined capital structure, dividends, and investment decisions.\\
These ideas influenced household economics, pension policy, and financial theory. They also provided a basis for studying how people plan consumption over their lifetimes.

\item[1986 --- James Buchanan]
James Buchanan founded public-choice analysis of political decision-making.\\
He treated voters, politicians, and bureaucrats as participants responding to incentives.\\
His constitutional economics studied rules that constrain collective choices and fiscal behavior.\\
He showed how institutional design can limit rent-seeking, deficits, and government failure. His work encouraged economists to study political rules as well as market rules.

\item[1987 --- Robert Solow]
Robert Solow developed the neoclassical model of economic growth.\\
He separated the contributions of capital, labor, and technological progress to output.\\
His growth accounting showed that long-run productivity improvements cannot be explained by inputs alone.\\
The Solow model remains a starting point for research on productivity and development. It also provides a standard way to separate capital accumulation from technological progress.

\item[1988 --- Maurice Allais]
Maurice Allais advanced the theory of competitive markets and efficient resource allocation.\\
He studied general equilibrium, capital accumulation, and the foundations of welfare economics.\\
His analysis of risk preferences produced the Allais paradox, challenging expected-utility assumptions.\\
His work combined mathematical theory with practical questions about markets and public policy. The Allais paradox also encouraged later research on decisions under risk.

\item[1989 --- Trygve Haavelmo]
Trygve Haavelmo established the statistical foundations for testing economic ideas.\\
He argued that economic relationships should be expressed as probabilistic models rather than exact identities.\\
His work clarified how data can be used to estimate connected relationships and identify causes.\\
These principles underpin modern causal and empirical economic research. They clarified how economic data can be used to test theories and estimate relationships.

\item[1990 --- Harry Markowitz, Merton Miller, and William Sharpe]
Harry Markowitz founded modern portfolio theory by showing how diversification can reduce risk for a given return.\\
Merton Miller developed the irrelevance propositions linking firm value to financing and dividend decisions under ideal conditions.\\
William Sharpe created the capital asset pricing model, relating expected returns to systematic risk.\\
Together they supplied core tools for asset pricing, corporate finance, and investment management. Their models are still used to measure risk and evaluate financial decisions.

\item[1991 --- Ronald Coase]
Ronald Coase showed that transaction costs determine whether economic activity is organized through markets or firms.\\
His theory of the firm explained why entrepreneurs coordinate transactions internally.\\
The Coase theorem clarified how property rights and bargaining can affect externalities when transaction costs are low.\\
His insights founded modern law and economics and institutional analysis. They also showed why property rights and enforcement costs matter for economic outcomes.

\item[1992 --- Gary Becker]
Gary Becker extended microeconomic reasoning to education, family decisions, crime, and discrimination.\\
He treated time, skills, and household choices as scarce resources allocated through incentives.\\
His human-capital theory explained how training and schooling affect earnings and productivity.\\
Becker's approach greatly expanded the range of behavior studied with economic models. It influenced research on education, households, crime, discrimination, and human capital.

\item[1993 --- Robert Fogel and Douglass North]
Robert Fogel used quantitative evidence and counterfactual analysis to reinterpret economic history.\\
He studied railroads, slavery, nutrition, and technological change with economic measurement.\\
Douglass North explained how institutions, property rights, and transaction costs shape long-run growth.\\
Together they made economic history more empirical, analytical, and institutionally focused. Their methods helped explain how past decisions affect present economic performance.

\item[1994 --- John Nash, Reinhard Selten, and John Harsanyi]
John Nash introduced equilibrium concepts for non-cooperative games and bargaining problems.\\
Reinhard Selten refined equilibrium theory through subgame perfection and analysis of sequential choices.\\
John Harsanyi developed methods for games with incomplete information and private types.\\
Their work made strategic interaction a central framework across economics and related social sciences. Game theory is now used to study competition, bargaining, voting, and international relations.

\item[1995 --- Robert Lucas]
Robert Lucas showed that people's expectations about future policies affect their decisions today.\\
He showed that people adapt their expectations to policy rules, changing the effects of predictable interventions.\\
He warned that a policy cannot be evaluated correctly if the analysis assumes that people will not change their behavior after the policy changes.\\
His work encouraged micro-founded, dynamic models of economic fluctuations. It also changed how economists evaluate the effects of government policy.

\item[1996 --- James Mirrlees and William Vickrey]
James Mirrlees studied how governments can design fair taxes when they cannot fully observe people's ability or effort.\\
His models explained the balance between reducing inequality and preserving incentives to work.\\
William Vickrey applied the study of incentives to auctions, road charges, and government finance.\\
Together they established foundations for the modern economics of incentives and mechanism design. Their ideas are applied to taxation, auctions, public services, and regulation.

\item[1997 --- Robert Merton and Myron Scholes]
Robert Merton and Myron Scholes developed a method for valuing stock options and similar financial contracts.\\
Their method connected investment prices, changes in risk, interest rates, and time.\\
It provided a practical way to hedge and manage the risks of contingent claims.\\
Their work transformed financial engineering and modern risk management. The pricing framework also provided a common language for trading and hedging derivatives.

\item[1998 --- Amartya Sen]
Amartya Sen broadened the study of well-being beyond income to include people's real opportunities and freedoms.\\
His study of poverty and famine showed that people can starve even when food exists if they cannot obtain it.\\
He developed measures and arguments concerning inequality, social choice, and human development.\\
His work placed individual opportunity and dignity at the heart of economic evaluation. It influenced poverty measurement and the development of the capability approach.

\item[1999 --- Robert Mundell]
Robert Mundell studied how money and government spending work under different systems for setting exchange rates.\\
His model showed how these policies depend on how easily money crosses borders and how exchange rates are set.\\
He developed the theory of optimum currency areas, informing debates about monetary unions.\\
His work strongly influenced international macroeconomics and the design of the euro. It also clarified the policy choices faced by countries with open financial markets.

\item[2000 --- James Heckman]
James Heckman developed methods for correcting errors caused when a study includes only selected people.\\
He showed how participation choices can distort estimates of wages and training effects.\\
His models separate observed results from the decision to take part in a program.\\
His methods improved labor studies and the evaluation of public programs. They remain important when researchers cannot observe all relevant choices directly.

\item[2000 --- Daniel McFadden]
Daniel McFadden developed statistical methods for studying choices between alternatives.\\
His models explain choices about transport, work, housing, and consumption.\\
He connected theories of choice with the analysis of individual data.\\
His work improved the measurement of decisions in applied economics. These models are useful when choices are recorded as separate alternatives.

\item[2001 --- George Akerlof, Michael Spence, and Joseph Stiglitz]
George Akerlof showed that markets can work poorly when buyers cannot tell good products from bad ones.\\
Michael Spence explained how education and other visible signs can give employers and buyers information about quality.\\
Joseph Stiglitz studied how lenders and other decision-makers respond when they do not have complete information.\\
Together they established asymmetric information as a central explanation of market outcomes. Their work explains why markets may need disclosure rules, contracts, or regulation.

\item[2002 --- Daniel Kahneman]
Daniel Kahneman used findings from psychology to study how people make economic decisions.\\
His prospect theory explained how people judge gains and losses.\\
He studied common errors in judgment and the effects of how choices are presented.\\
His work helped establish behavioral economics. It showed that psychological evidence can improve economic models of individual choice.

\item[2002 --- Vernon Smith]
Vernon Smith developed laboratory experiments to study markets and rules.\\
His experiments showed how prices and trade outcomes change under different market designs.\\
He created methods for comparing economic theory with observed behavior.\\
His work helped establish experimental economics. Laboratory evidence remains useful for testing market rules before practical use.

\item[2003 --- Robert Engle and Clive Granger]
Robert Engle developed ARCH models for time series whose volatility changes over time.\\
His methods allow economists to estimate and forecast changing financial risk.\\
Clive Granger developed the concept of cointegration for non-stationary economic variables with stable long-run relationships.\\
Together they transformed empirical analysis of financial and macroeconomic time series. Their methods remain standard tools for forecasting and studying long-run relationships.

\item[2004 --- Finn Kydland and Edward Prescott]
Finn Kydland and Edward Prescott showed how time inconsistency can weaken discretionary economic policy.\\
Their work explained why credible rules may outperform policies repeatedly chosen after conditions change.\\
They also developed real-business-cycle analysis, emphasizing technology and intertemporal optimization.\\
Their research shaped dynamic macroeconomics and the study of policy credibility. It also provided a framework for analyzing how expectations affect economic cycles.

\item[2005 --- Robert Aumann and Thomas Schelling]
Robert Aumann advanced repeated-game theory and explained how cooperation can persist among self-interested players.\\
His work also clarified the role of common knowledge in strategic reasoning.\\
Thomas Schelling analyzed bargaining, conflict, deterrence, and coordination in social interactions.\\
Together they showed how strategic incentives can produce both conflict and stable cooperation. Their ideas are applied to bargaining, security, and repeated business relationships.

\item[2006 --- Edmund Phelps]
Edmund Phelps analyzed the intertemporal trade-offs faced by macroeconomic policy.\\
He distinguished short-run inflation--employment effects from the long-run natural rate of unemployment.\\
His work emphasized expectations, wage-setting, and the role of incentives in employment.\\
These insights helped explain stagflation and shaped modern monetary policy analysis. They also emphasized the importance of expectations in employment and inflation.

\item[2007 --- Leonid Hurwicz, Eric Maskin, and Roger Myerson]
Leonid Hurwicz founded mechanism design, studying how rules can achieve desired outcomes with private information.\\
Eric Maskin developed implementation theory, identifying when social goals can be enforced through mechanisms.\\
Roger Myerson advanced optimal auction design and incentive-compatible mechanisms.\\
Their work provides the theoretical basis for designing markets, voting systems, and public institutions. It explains how rules can be chosen when participants have different information and interests.

\item[2008 --- Paul Krugman]
Paul Krugman explained trade patterns that arise from economies of scale and imperfect competition.\\
His new trade theory showed why similar countries may trade heavily with one another.\\
His work on economic geography linked market size, transport costs, and the location of production.\\
These ideas connected international trade with regional development and urban concentration. They also explained why scale and location influence the pattern of world production.

\item[2009 --- Elinor Ostrom]
Elinor Ostrom showed that communities can manage shared resources through local rules.\\
Her research showed that common resources do not always require private or central control.\\
She identified the importance of monitoring, trust, sanctions, and cooperation.\\
Her work changed the study of shared resources and local institutions. It showed the value of studying real rules and outcomes as well as formal models.

\item[2009 --- Oliver Williamson]
Oliver Williamson explained why some activities take place inside firms rather than in markets.\\
He showed how the cost of making and enforcing agreements affects this choice.\\
He studied contracts, firm organization, and the limits of market exchange.\\
His research established transaction-cost economics. The approach remains useful for studying firms, contracts, and economic institutions.

\item[2010 --- Peter Diamond, Dale Mortensen, and Christopher Pissarides]
Peter Diamond analyzed search costs and how they affect prices, unemployment, and public policy.\\
Dale Mortensen modeled job search, matching, worker turnover, and equilibrium unemployment.\\
Christopher Pissarides developed matching models that explain how vacancies and job seekers meet.\\
Their framework made labor-market frictions central to modern macroeconomics. It is used to study unemployment, job vacancies, wages, and labor-market policy.

\item[2011 --- Thomas Sargent and Christopher Sims]
Thomas Sargent used structural macroeconometric methods to study expectations and policy changes.\\
His work clarified how economic agents learn and how policy regimes affect outcomes.\\
Christopher Sims developed vector autoregression and innovations accounting for identifying macroeconomic shocks.\\
Together they supplied influential empirical tools for analyzing cause and effect in the macroeconomy. Their methods help distinguish economic shocks from simple correlations.

\item[2012 --- Alvin Roth and Lloyd Shapley]
Lloyd Shapley developed mathematical theories of stable matching and cooperative allocation.\\
His deferred-acceptance ideas explain how participants can be matched without incentives for destabilizing changes.\\
Alvin Roth applied these principles to school choice, labor markets, and organ exchange.\\
Together they founded modern market design, turning matching theory into practical institutions. Their methods are used in school admissions, medical matching, and labor markets.

\item[2013 --- Eugene Fama, Lars Hansen, and Robert Shiller]
Eugene Fama developed tests of asset-price efficiency and the predictability of returns.\\
Lars Hansen created the generalized method of moments for estimating dynamic asset-pricing models.\\
Robert Shiller documented long-run excess volatility and limited predictability in asset markets.\\
Their contrasting findings advanced empirical research on how financial prices are formed. They showed both the strength of market competition and the limits of simple efficiency claims.

\item[2014 --- Jean Tirole]
Jean Tirole developed modern theories of market power, strategic competition, and regulation.\\
He analyzed how information asymmetries affect firms, regulators, and consumers.\\
His models address oligopoly, natural monopolies, banking, innovation, and digital-platform markets.\\
His work gives policymakers tools for designing regulation in complex industries. It is applied to competition policy, telecommunications, banking, and digital markets.

\item[2015 --- Angus Deaton]
Angus Deaton developed rigorous methods for studying household consumption and demand.\\
He linked individual consumption choices with aggregate saving and economic growth.\\
His research measured poverty and welfare while emphasizing differences among households.\\
These methods improved development policy and the use of household survey data. They also support comparisons of living standards across different households.

\item[2016 --- Oliver Hart and Bengt Holmström]
Oliver Hart developed incomplete-contract theory for situations in which future contingencies cannot be fully specified.\\
His work explained how ownership and control should be allocated within firms and organizations.\\
Bengt Holmström analyzed incentives, moral hazard, contracts, and optimal compensation.\\
Together they established modern contract theory for firms, finance, and public institutions. Their work explains how contracts can share risk and provide incentives when information is incomplete.

\item[2017 --- Richard Thaler]
Richard Thaler showed how limited rationality and self-control affect saving, spending, and market behavior.\\
He documented mental accounting, endowment effects, and other systematic departures from standard models.\\
His research developed the idea of choice architecture and gentle ``nudges.''\\
He made behavioral insights useful for policy, finance, and everyday economic decisions. His work also encouraged policies that help people without removing their freedom to choose.

\item[2018 --- William Nordhaus]
William Nordhaus integrated climate change into long-run economic growth and policy analysis.\\
His models evaluated the costs of emissions and the benefits of climate policy.\\
He showed how carbon pricing and other policies can address the global externality of emissions.\\
His work established climate economics as a central part of long-run policy analysis. His models help compare the present costs of climate action with future benefits.

\item[2018 --- Paul Romer]
Paul Romer developed endogenous-growth theory, showing how ideas and innovation can generate persistent growth.\\
His models treated knowledge as a non-rival input whose benefits can extend across the economy.\\
He explained how research, human capital, and policy incentives influence technological progress.\\
His work placed innovation and ideas at the center of growth theory. The framework explains why research and knowledge can support continuing economic growth.

\item[2019 --- Abhijit Banerjee, Esther Duflo, and Michael Kremer]
Abhijit Banerjee and Esther Duflo pioneered field experiments that test development policies in real communities.\\
Their work used carefully designed interventions to identify effective responses to poverty.\\
Michael Kremer applied randomized experiments to education, health, and agricultural development.\\
Together they made evidence-based, problem-focused evaluation central to modern development economics. Their approach helps policymakers compare programs using measured effects.

\item[2020 --- Paul Milgrom and Robert Wilson]
Robert Wilson analyzed auctions in which bidders have interdependent values and uncertain information.\\
Paul Milgrom extended auction theory and studied how information affects bidding and market outcomes.\\
They designed practical formats, including spectrum auctions, that allocate complex resources efficiently.\\
Their work improved both the theory and practice of market design. Their auction methods are used to allocate public resources such as radio spectrum.

\item[2021 --- David Card]
David Card used natural experiments to study minimum wages, immigration, education, and labor-market outcomes.\\
His evidence challenged simple assumptions about wages, employment, and the effects of immigration.\\
His work demonstrated how natural experiments can answer important labor-economics questions. It encouraged wider use of institutional changes as sources of evidence.

\item[2021 --- Joshua Angrist]
Joshua Angrist developed instrumental-variable methods for estimating causal effects with imperfect treatment assignment.\\
He applied quasi-experimental designs to education, schooling, military service, and labor markets.\\
His research clarified how researchers can use external variation to identify treatment effects.\\
His work strengthened the empirical foundations of modern causal inference. It provides methods for estimating effects when controlled experiments are not possible.

\item[2021 --- Guido Imbens]
Guido Imbens clarified identification, treatment effects, and the interpretation of quasi-experimental evidence.\\
He established when instrumental variables and other research designs support causal conclusions.\\
His work distinguished local effects, heterogeneous effects, and the populations to which estimates apply.\\
His methods made empirical causal claims more transparent and reliable. They require researchers to state which effects are identified and for whom they apply.

\item[2022 --- Ben Bernanke, Douglas Diamond, and Philip Dybvig]
Ben Bernanke showed how bank failures and credit disruptions can deepen economic depressions.\\
Douglas Diamond explained the role of banks in monitoring borrowers and transforming maturities.\\
Philip Dybvig modeled bank runs and showed why deposit insurance can stabilize financial systems.\\
Together they provided essential foundations for understanding banking and financial crises. Their work informs deposit insurance, bank supervision, and crisis policy.

\item[2023 --- Claudia Goldin]
Claudia Goldin reconstructed the long history of women's participation and wages in the labor market.\\
She identified how education, technology, marriage, parenthood, and occupations shaped gender gaps.\\
Her research showed why progress toward equality has been uneven across generations.\\
Her historical evidence informs current debates about work, care, and pay inequality. It also shows how labor-market institutions affect social change.

\item[2024 --- Daron Acemoglu, Simon Johnson, and James Robinson]
Daron Acemoglu, Simon Johnson, and James Robinson showed how political and economic institutions shape prosperity.\\
They distinguished inclusive institutions, which broaden opportunity, from extractive institutions, which concentrate power.\\
Their historical research traced how colonial origins and critical junctures created persistent development paths.\\
Their work placed governance, incentives, and institutional change at the center of development economics. The framework is used to study why development paths differ across countries and over time.

\end{description}}

\end{document}
